\documentclass[aspectratio=169,xcolor=dvipsnames]{beamer}
\usetheme{Simple}
\usepackage[T1]{fontenc}
\usepackage[utf8]{inputenc}
\usepackage[polish]{babel}
\usepackage{graphicx}
\usepackage{booktabs}
\usepackage{tabularx}
\newcolumntype{Y}{>{\raggedright\arraybackslash}X}
\usepackage{xcolor}
\usepackage{hyperref}
\usepackage{listings}
\usepackage{tikz}
\usetikzlibrary{positioning,arrows.meta,fit,backgrounds}

\definecolor{CyberBlue}{HTML}{2464AD}
\definecolor{CyberNavy}{HTML}{173D70}
\definecolor{CyberPale}{HTML}{E8F2FC}
\definecolor{CyberText}{HTML}{292525}
\definecolor{CyberRed}{HTML}{B23A3A}
\definecolor{CyberGreen}{HTML}{2E7D5B}

% Nadpisanie barw motywu Simple
\setbeamercolor{normal text}{fg=CyberText,bg=white}
\setbeamercolor{structure}{fg=CyberBlue}
\setbeamercolor{title}{fg=white}
\setbeamercolor{subtitle}{fg=white}
\setbeamercolor{frametitle}{fg=CyberNavy,bg=white}
\setbeamercolor{section in toc}{fg=CyberNavy}
\setbeamercolor{item}{fg=CyberBlue}
\setbeamercolor{subitem}{fg=CyberBlue}
\setbeamercolor{enumerate item}{fg=CyberBlue}
\setbeamercolor{alerted text}{fg=CyberBlue}
\setbeamercolor{block title}{fg=white,bg=CyberNavy}
\setbeamercolor{block body}{fg=CyberText,bg=CyberPale}
\setbeamercolor{block title alerted}{fg=white,bg=CyberBlue}
\setbeamercolor{block body alerted}{fg=CyberText,bg=CyberPale}
\setbeamercolor{block title example}{fg=white,bg=CyberNavy}
\setbeamercolor{block body example}{fg=CyberText,bg=CyberPale}
\setbeamercolor{palette primary}{fg=white,bg=CyberBlue}
\setbeamercolor{palette secondary}{fg=white,bg=CyberNavy}
\setbeamercolor{palette tertiary}{fg=white,bg=CyberNavy}
\setbeamercolor{footline}{fg=CyberNavy,bg=white}
\setbeamercolor{navigation symbols}{fg=CyberBlue}
\hypersetup{colorlinks=true,urlcolor=CyberBlue,linkcolor=CyberNavy,citecolor=CyberBlue}

% Listingi (kod bez polskich znaków)
\lstdefinestyle{cyber}{
  basicstyle=\ttfamily\scriptsize,
  keywordstyle=\color{CyberBlue}\bfseries,
  commentstyle=\color{gray}\itshape,
  stringstyle=\color{CyberGreen},
  backgroundcolor=\color{CyberPale},
  frame=single, rulecolor=\color{CyberBlue!40},
  framesep=3pt, xleftmargin=4pt, xrightmargin=4pt,
  showstringspaces=false, columns=fullflexible, keepspaces=true,
  breaklines=true
}
\lstset{style=cyber}

% Znacznik miejsc do uzupełnienia po testach
\newcommand{\todo}[1]{\textcolor{CyberRed}{\textbf{[TODO:} #1\textbf{]}}}

% Podpis slajdu: kto omawia
\newcommand{\kto}[1]{%
  \begin{tikzpicture}[remember picture,overlay]
    \node[anchor=south west, fill=CyberPale, draw=CyberBlue!50, rounded corners=2pt,
          font=\tiny, text=CyberNavy, inner sep=2pt]
      at ([xshift=4mm,yshift=1mm]current page.south west)
      {\textbf{#1}};
  \end{tikzpicture}}

% Diagramy
\tikzset{
  box/.style={draw=CyberNavy, fill=CyberPale, rounded corners=3pt, thick,
              align=center, font=\scriptsize, inner sep=4pt, minimum height=0.9cm},
  hot/.style={box, draw=CyberRed, fill=CyberRed!8},
  arr/.style={-{Stealth[length=2mm]}, thick, CyberNavy},
  lbl/.style={font=\tiny, text=CyberText, align=center}
}

\title[CrowdSec -- ochrona WordPressa]{Ochrona aplikacji webowej przed atakami HTTP\\z użyciem CrowdSec (IPS + WAF AppSec)}
\subtitle{Cyberbezpieczeństwo -- Moduł 2: Ochrona}
\author[D. Pągowski, OSOBA 2, OSOBA 3]{Dawid Pągowski \and OSOBA 2 \and OSOBA 3}
\institute[PŚ]{Wydział Matematyki Stosowanej\\Politechnika Śląska}
\date{\today}

\begin{document}

\begin{frame}
    \titlepage
\end{frame}

\begin{frame}{Plan prezentacji}
    \tableofcontents
\end{frame}

% =====================================================================
\section{Wprowadzenie}
% =====================================================================

\begin{frame}{Problem: zagrożenie i chroniony zasób}
    \kto{Dawid Pągowski}
    \begin{columns}[T]
        \column{0.5\textwidth}
        \begin{block}{Chroniony zasób}
            \textbf{WordPress 6} za nginx (reverse proxy) + MariaDB.
        \end{block}
        \begin{block}{Scenariusze zagrożeń}
            \begin{itemize}\small
                \item skanowanie i rekonesans (pliki wrażliwe, panele admina),
                \item brute force logowania (\texttt{wp-login.php}),
                \item ataki na treść żądań: SQLi, XSS, path traversal,
                \item exploity znanych CVE.
            \end{itemize}
        \end{block}
        \column{0.46\textwidth}
        \begin{exampleblock}{Kryteria oceny}
            \begin{itemize}\small
                \item \textbf{wykrycie} -- alert w CrowdSec,
                \item \textbf{zablokowanie} -- odpowiedź 403 / ban IP,
                \item \textbf{skuteczność ataku} w wariancie bez i z ochroną,
                \item \textbf{czas / liczba żądań} do bana,
                \item \textbf{brak false positive} dla zwykłego ruchu.
            \end{itemize}
        \end{exampleblock}
    \end{columns}
\end{frame}

\begin{frame}{Podział pracy w zespole}
    \begin{columns}[T]
        \column{0.32\textwidth}
        \begin{block}{Dawid Pągowski}
            \small
            \begin{itemize}
                \item przegląd rozwiązań i wybór CrowdSec,
                \item środowisko \texttt{docker compose},
                \item konfiguracja CrowdSec, bouncera i AppSec,
                \item problemy i ograniczenia.
            \end{itemize}
        \end{block}
        \column{0.32\textwidth}
        \begin{block}{OSOBA 2}
            \small
            \begin{itemize}
                \item \textbf{atak: SQL injection} -- własny skrypt,
                \item test z wyłączonym i włączonym CrowdSec.
            \end{itemize}
        \end{block}
        \column{0.32\textwidth}
        \begin{block}{OSOBA 3}
            \small
            \begin{itemize}
                \item \textbf{atak: brute force logowania WordPress} -- własny skrypt,
                \item test z wyłączonym i włączonym CrowdSec.
            \end{itemize}
        \end{block}
    \end{columns}
\end{frame}

% =====================================================================
\section{Przegląd rozwiązań i wybór}
% =====================================================================

\begin{frame}{Przegląd dostępnych rozwiązań}
    \kto{Dawid Pągowski}
    \scriptsize
    \renewcommand{\arraystretch}{1.15}
    \begin{tabularx}{\textwidth}{@{}>{\raggedright\arraybackslash}p{2.3cm} Y c c l Y@{}}
        \toprule
        \textbf{Rozwiązanie} & \textbf{Zasada działania} & \textbf{WAF} & \textbf{Reputacja IP} & \textbf{Licencja} & \textbf{Główne ograniczenie} \\
        \midrule
        fail2ban & wyrażenia regularne na logach $\rightarrow$ ban w zaporze (iptables/nftables) & -- & -- & GPL-2.0 & nie widzi treści żądań, działa po fakcie, lokalny \\
        ModSecurity / Coraza + OWASP CRS & reguły inspekcji każdego żądania (WAF); Coraza -- następca w Go & \checkmark & -- & Apache-2.0 & brak analizy behawioralnej, wymaga tuningu (false positive) \\
        \textbf{CrowdSec} & scenariusze na logach + AppSec (WAF) + bouncery & \checkmark & \checkmark & MIT (silnik) & fail-open, część funkcji w płatnej konsoli \\
        WAF chmurowy (np. Cloudflare) & proxy u dostawcy & \checkmark & \checkmark & płatny / freemium & nie da się przetestować w izolowanym labie \\
        \bottomrule
    \end{tabularx}
\end{frame}

\begin{frame}{fail2ban vs CrowdSec}
    \kto{Dawid Pągowski}
    \begin{columns}[T]
        \column{0.48\textwidth}
        \begin{block}{fail2ban}
            \small
            \begin{itemize}
                \item \emph{jail} = filtr (regex) + akcja (ban w firewallu),
                \item prosty, sprawdzony, minimalne zasoby,
                \item jedna maszyna -- brak współdzielenia wiedzy,
                \item tylko analiza logów: każdy atak musi najpierw \emph{dotrzeć} do aplikacji,
                \item brak inspekcji treści żądań -- pojedynczy udany SQLi nie zostanie zatrzymany.
            \end{itemize}
        \end{block}
        \column{0.48\textwidth}
        \begin{block}{CrowdSec}
            \small
            \begin{itemize}
                \item parsery + scenariusze (\emph{leaky bucket}) zamiast regexów,
                \item gotowe kolekcje z Hub (nginx, WordPress, CVE),
                \item rozdzielenie detekcji (agent/LAPI) od egzekucji (bouncer),
                \item \textbf{AppSec} -- blokada \emph{w trakcie} żądania (virtual patching, reguły generyczne),
                \item blocklisty społeczności (reputacja IP).
            \end{itemize}
        \end{block}
    \end{columns}
\end{frame}

\begin{frame}{Uzasadnienie wyboru: CrowdSec}
    \kto{Dawid Pągowski}
    \begin{block}{Dlaczego CrowdSec}
        \begin{itemize}
            \item jako jedyne darmowe rozwiązanie łączy \textbf{detekcję behawioralną} (jak fail2ban) z \textbf{inspekcją żądań} (jak WAF) -- dwie warstwy obrony w jednym narzędziu,
            \item gotowe scenariusze dla nginx i WordPressa -- mniej własnych reguł do utrzymania,
            \item bouncer w nginx blokuje na poziomie HTTP (403), a nie tylko IP w zaporze,
            \item działa w kontenerach -- powtarzalne, izolowane środowisko laboratoryjne,
            \item licencja MIT, aktywna dokumentacja i społeczność.
        \end{itemize}
    \end{block}
    \begin{alertblock}{Czego się spodziewamy}
        Pozwala porównać obie warstwy na tym samym ruchu: AppSec (pojedyncze żądanie) vs scenariusze (seria zdarzeń).
    \end{alertblock}
\end{frame}

% =====================================================================
\section{Charakterystyka rozwiązania ochronnego}
% =====================================================================

\begin{frame}{Zasada działania CrowdSec}
    \kto{Dawid Pągowski}
    \centering
    \begin{tikzpicture}[node distance=0.45cm and 0.45cm]
        % Ścieżka detekcji z logów
        \node[box] (log) {Logi\\nginx};
        \node[box, right=of log] (parse) {Parsery\\(s00--s02)};
        \node[box, right=of parse] (scen) {Scenariusze\\\emph{leaky bucket}};
        \node[box, right=0.7cm of scen] (lapi) {LAPI\\decyzja: ban};
        \node[box, right=of lapi] (bouncer) {Bouncer\\(nginx + Lua)};
        \draw[arr] (log) -- (parse);
        \draw[arr] (parse) -- (scen);
        \draw[arr] (scen) -- node[lbl, below=1pt]{alert} (lapi);
        \draw[arr] (bouncer) to[bend right=35] node[lbl, above]{czy IP ma bana?} (lapi);
        % Ścieżka AppSec
        \node[hot, below=1.3cm of bouncer] (appsec) {AppSec\\(reguły WAF)};
        \node[box, left=of appsec, xshift=-1.3cm] (resp) {403 albo przekazanie\\do aplikacji};
        \draw[arr, dashed] (bouncer) -- node[lbl, right]{każde żądanie\\(in-band)} (appsec);
        \draw[arr] (appsec) -- node[lbl, above]{werdykt} (resp);
        % Opisy warstw
        \node[lbl, above=0.15cm of parse, xshift=1.1cm, font=\scriptsize\bfseries, text=CyberNavy] {detekcja behawioralna (po fakcie)};
        \node[lbl, below=0.15cm of resp, xshift=1.6cm, font=\scriptsize\bfseries, text=CyberRed] {inspekcja żądania (zanim dotrze do aplikacji)};
    \end{tikzpicture}
    \vspace{0.2cm}

    \small Detekcja (agent, LAPI) i egzekucja (bouncer) są rozdzielone -- jeden agent może zasilać wiele bouncerów.
\end{frame}

\begin{frame}{Rozwiązanie, ograniczenia, zastosowanie}
    \kto{Dawid Pągowski}
    \begin{block}{Rozwiązanie}
        \small CrowdSec 1.x w kontenerze: kolekcje \texttt{crowdsecurity/nginx}, \texttt{crowdsecurity/wordpress}, \texttt{appsec-virtual-patching}, \texttt{appsec-generic-rules}; bouncer \texttt{crowdsec-nginx-bouncer} (moduł Lua) z włączonym AppSec.
    \end{block}
    \begin{alertblock}{Ograniczenia}
        \small Scenariusze reagują dopiero po serii zdarzeń; bouncer działa \emph{fail-open}; nie chroni przed błędami logiki aplikacji ani logowaniem poprawnym (wykradzionym) hasłem.
    \end{alertblock}
    \begin{exampleblock}{Przykład zastosowania}
        \small Publiczny WordPress: AppSec odcina znane exploity od pierwszego żądania, scenariusze banują skanery i brute force.
    \end{exampleblock}
\end{frame}

\begin{frame}{Architektura środowiska testowego}
    \kto{Dawid Pągowski}
    \centering
    \begin{tikzpicture}[node distance=0.55cm and 0.9cm]
        \node[hot] (att) {\textbf{attacker}\\nicolaka/netshoot\\(narzędzia ataku)};
        \node[box, right=of att] (ngx) {\textbf{nginx}\\Debian + Lua\\bouncer, :80};
        \node[box, right=of ngx] (wp) {\textbf{wordpress}\\6-apache};
        \node[box, right=of wp] (db) {\textbf{db}\\MariaDB 11};
        \node[box, below=1.0cm of ngx, minimum width=3.2cm] (cs) {\textbf{crowdsec}\\agent + LAPI :8080\\AppSec :7422};
        \draw[arr] (att) -- node[lbl, above]{HTTP} (ngx);
        \draw[arr] (ngx) -- (wp);
        \draw[arr] (wp) -- (db);
        \draw[arr] ([xshift=-4mm]ngx.south) -- node[lbl, left]{access.log\\(wolumen)} ([xshift=-4mm]cs.north);
        \draw[arr, {Stealth[length=2mm]}-{Stealth[length=2mm]}] ([xshift=4mm]ngx.south) -- node[lbl, right]{decyzje,\\AppSec} ([xshift=4mm]cs.north);
        \begin{scope}[on background layer]
            \node[draw=CyberBlue, dashed, rounded corners, fit=(att)(ngx)(wp)(db)(cs), inner sep=8pt,
                  label={[font=\tiny, text=CyberBlue]below right:sieć Docker \texttt{internal} 10.11.0.0/16}] {};
        \end{scope}
    \end{tikzpicture}
    \vspace{0.15cm}

    \small Całość w jednym \texttt{docker compose} -- izolowane środowisko, ataki wyłącznie na własne kontenery.
\end{frame}

\begin{frame}[fragile]{Konfiguracja: usługa CrowdSec w \texttt{docker compose}}
    \kto{Dawid Pągowski}
    \begin{columns}[T]
        \column{0.55\textwidth}
\begin{lstlisting}[language=bash]
crowdsec:
  image: crowdsecurity/crowdsec
  environment:
    COLLECTIONS: "crowdsecurity/nginx
      crowdsecurity/wordpress
      crowdsecurity/appsec-virtual-patching
      crowdsecurity/appsec-generic-rules"
    DISABLE_PARSERS: "crowdsecurity/whitelists"
    BOUNCER_KEY_nginx: ${BOUNCER_KEY}
  volumes:
    - ./crowdsec/acquis.yaml:/etc/crowdsec/acquis.yaml
    - ./crowdsec/acquis.d:/etc/crowdsec/acquis.d
    - ./nginx/logs:/var/log/nginx
  healthcheck:
    test: ["CMD", "cscli", "lapi", "status"]
\end{lstlisting}
        \column{0.42\textwidth}
        \small
        \begin{itemize}
            \item kolekcje z Hub instalowane przy starcie kontenera,
            \item klucz bouncera rejestrowany automatycznie (\texttt{BOUNCER\_KEY\_nginx}),
            \item logi nginx współdzielone przez wolumen,
            \item \texttt{acquis.d} -- źródło logów + komponent AppSec (:7422),
            \item healthcheck bazy (\texttt{healthcheck.sh --innodb\_initialized}) -- WordPress startuje dopiero po gotowości MariaDB.
        \end{itemize}
    \end{columns}
\end{frame}

\begin{frame}[fragile]{Konfiguracja: nginx z bouncerem}
    \kto{Dawid Pągowski}
    \begin{columns}[T]
        \column{0.55\textwidth}
\begin{lstlisting}[language=bash]
FROM debian:bookworm-slim
RUN apt-get install -y nginx \
      libnginx-mod-http-lua lua-cjson \
      crowdsec-nginx-bouncer
RUN sed -i 's|^APPSEC_URL=.*|APPSEC_URL=http://crowdsec:7422|' \
    /etc/crowdsec/bouncers/crowdsec-nginx-bouncer.conf
COPY docker_resolver.conf \
     /etc/nginx/conf.d/00_docker_resolver.conf
\end{lstlisting}
\begin{lstlisting}
# 00_docker_resolver.conf
resolver 127.0.0.11 valid=30s ipv6=off;
\end{lstlisting}
        \column{0.42\textwidth}
        \small
        \begin{itemize}
            \item bouncer działa jako moduł Lua w fazie \emph{access} -- każde żądanie sprawdzane przed przekazaniem do WordPressa,
            \item adres LAPI i klucz przekazywane zmiennymi środowiskowymi,
            \item AppSec wskazany na kontener \texttt{crowdsec},
            \item \texttt{resolver} wymagany, bo Lua w nginx nie korzysta z \texttt{/etc/resolv.conf}.
        \end{itemize}
    \end{columns}
\end{frame}

\begin{frame}{Problemy napotkane podczas konfiguracji}
    \kto{Dawid Pągowski}
    \footnotesize
    \renewcommand{\arraystretch}{1.2}
    \begin{tabularx}{\textwidth}{@{}>{\raggedright\arraybackslash}p{3.6cm} Y Y@{}}
        \toprule
        \textbf{Objaw} & \textbf{Przyczyna} & \textbf{Rozwiązanie} \\
        \midrule
        \texttt{no resolver defined to resolve "crowdsec"} & cosockety Lua używają tylko dyrektywy \texttt{resolver} nginx & \texttt{resolver 127.0.0.11} (DNS Dockera) \\
        \texttt{crowdsec could not be resolved (3: Host not found)} & kontenery w różnych sieciach Dockera & wspólna sieć \texttt{internal} \\
        ataki z sieci 10.11.0.0/16 ignorowane & domyślny parser \texttt{crowdsecurity/whitelists} pomija adresy prywatne & \texttt{DISABLE\_PARSERS} \\
        WordPress startuje przed bazą & MariaDB inicjalizuje się ok. 1 min & healthcheck + \texttt{condition: service\_healthy} \\
        \bottomrule
    \end{tabularx}
    \vspace{0.1cm}

    \footnotesize Przy dwóch pierwszych błędach bouncer \textbf{przepuszczał cały ruch}.
\end{frame}

% =====================================================================
\section{Praca laboratoryjna i wyniki}
% =====================================================================

\begin{frame}{Realizacja eksperymentu}
    \kto{Dawid Pągowski}
    \begin{columns}[T]
        \column{0.5\textwidth}
        \begin{block}{Metoda}
            \small
            \begin{enumerate}
                \item Każdy atak przygotowuje i wykonuje jedna osoba (własny skrypt).
                \item \textbf{Wariant A} -- CrowdSec wyłączony (\texttt{docker compose stop crowdsec}).
                \item \textbf{Wariant B} -- CrowdSec włączony.
                \item Ataki z kontenera \texttt{attacker} w sieci laboratorium.
                \item Dowody: kody HTTP, \texttt{cscli alerts list}, \texttt{cscli decisions list}.
            \end{enumerate}
        \end{block}
        \column{0.46\textwidth}
        \begin{alertblock}{Powtarzalność testów}
            \small
            \begin{itemize}
                \item przed testem: \texttt{cscli decisions delete --all},
                \item scenariusze mają \emph{blackhole} (kilka minut ignorowania IP po alercie) -- między powtórzeniami restart CrowdSec,
                \item test zwykłego ruchu (przeglądanie, logowanie admina) -- false positive.
            \end{itemize}
        \end{alertblock}
    \end{columns}
\end{frame}

\begin{frame}{Zaobserwowane ograniczenia}
    \kto{Dawid Pągowski}
    \begin{itemize}
        \item \textbf{Fail-open} -- gdy bouncer nie może połączyć się z LAPI/AppSec, przepuszcza ruch; błąd konfiguracji DNS wyłączył całą ochronę bez widocznego objawu dla użytkownika.
        \item \textbf{Detekcja po fakcie} -- scenariusz przepuszcza pierwsze $N$ żądań serii ($N$ = pojemność kubełka); żądania te docierają do aplikacji.
        \item \textbf{Blackhole} -- po alercie dany scenariusz ignoruje IP przez kilka minut; po ręcznym zdjęciu bana ataki przechodzą.
        \item \textbf{Zaufanie do sieci prywatnych} -- domyślnie atakujący z wnętrza sieci (np. przejęty host) jest pomijany.
    \end{itemize}
\end{frame}


% ---------------- SQL injection (OSOBA 2) ----------------
\begin{frame}{Atak 1: SQL injection -- przygotowanie}
    \kto{OSOBA 2}
    \begin{columns}[T]
        \column{0.48\textwidth}
        \begin{block}{Na czym polega atak}
            \small Wstrzyknięcie fragmentu zapytania SQL w parametr lub pole formularza, aby odczytać lub zmienić dane w bazie.
        \end{block}
        \begin{block}{Która warstwa CrowdSec powinna zadziałać}
            \small \textbf{AppSec} (reguły WAF) -- blokada już pierwszego żądania kodem 403.
        \end{block}
        \column{0.48\textwidth}
        \begin{exampleblock}{Skrypt testowy}
            \small
            \todo{fragment własnego skryptu: adres celu, payloady, ile żądań}\\[4pt]
            \todo{kryterium: co uznajemy za udany atak, a co za blokadę}
        \end{exampleblock}
    \end{columns}
\end{frame}

\begin{frame}{Atak 1: SQL injection -- wyniki}
    \kto{OSOBA 2}
    \small
    \renewcommand{\arraystretch}{1.3}
    \begin{tabularx}{\textwidth}{@{}>{\raggedright\arraybackslash}p{3.6cm} Y Y@{}}
        \toprule
        & \textbf{CrowdSec wyłączony} & \textbf{CrowdSec włączony} \\
        \midrule
        Kody odpowiedzi HTTP & \todo{} & \todo{} \\
        Czy atak dotarł do aplikacji & \todo{} & \todo{} \\
        Alert / reguła w CrowdSec & -- & \todo{nazwa reguły z \texttt{cscli alerts list}} \\
        Ban adresu IP & -- & \todo{} \\
        \bottomrule
    \end{tabularx}
    \vspace{0.3cm}

    \todo{zrzuty z wynikami poleceń: \texttt{cscli alerts list}, \texttt{cscli alerts inspect <ID>}, \texttt{cscli decisions list}, \texttt{cscli metrics} oraz kodów HTTP ze skryptu -- dla obu wariantów; jedno zdanie wniosku}
\end{frame}

% ---------------- Brute force (OSOBA 3) ----------------
\begin{frame}{Atak 2: brute force logowania WordPress -- przygotowanie}
    \kto{OSOBA 3}
    \begin{columns}[T]
        \column{0.48\textwidth}
        \begin{block}{Na czym polega atak}
            \small Wielokrotne próby logowania do \texttt{wp-login.php} z różnymi hasłami, aż do trafienia poprawnego.
        \end{block}
        \begin{block}{Która warstwa CrowdSec powinna zadziałać}
            \small \textbf{Scenariusz z logów} (\texttt{crowdsecurity/http-bf-wordpress\_bf}) -- ban IP po serii nieudanych prób.
        \end{block}
        \column{0.48\textwidth}
        \begin{exampleblock}{Skrypt testowy}
            \small
            \todo{fragment własnego skryptu: login, lista haseł, odstęp między próbami}\\[4pt]
            \todo{kryterium: po ilu próbach następuje ban}
        \end{exampleblock}
    \end{columns}
\end{frame}

\begin{frame}{Atak 2: brute force logowania WordPress -- wyniki}
    \kto{OSOBA 3}
    \small
    \renewcommand{\arraystretch}{1.3}
    \begin{tabularx}{\textwidth}{@{}>{\raggedright\arraybackslash}p{3.6cm} Y Y@{}}
        \toprule
        & \textbf{CrowdSec wyłączony} & \textbf{CrowdSec włączony} \\
        \midrule
        Liczba wykonanych prób & \todo{} & \todo{} \\
        Próba, po której nastąpił ban & -- & \todo{} \\
        Czas do bana & -- & \todo{} \\
        Alert w CrowdSec & -- & \todo{nazwa scenariusza} \\
        Wolne próby (np. co 30 s) & -- & \todo{czy wykryto?} \\
        \bottomrule
    \end{tabularx}
    \vspace{0.3cm}

    \todo{zrzuty z wynikami poleceń: \texttt{cscli alerts list}, \texttt{cscli alerts inspect <ID>}, \texttt{cscli decisions list}, \texttt{cscli metrics} oraz kodów HTTP ze skryptu -- dla obu wariantów; jedno zdanie wniosku}
\end{frame}

% =====================================================================
\section{Wnioski}
% =====================================================================

\begin{frame}{Wnioski zespołu}
    \kto{OSOBA 3}
    \begin{columns}[T]
        \column{0.5\textwidth}
        \begin{block}{Ocena}
            \small
            \begin{itemize}
                \item CrowdSec ogranicza rekonesans, brute force i znane exploity HTTP,
                \item AppSec i scenariusze uzupełniają się: pojedyncze żądanie vs seria zdarzeń,
                \item jeden pozytywny test nie oznacza pełnej ochrony -- skuteczność zależy od konfiguracji i progów,
                \item \todo{potwierdzić wynikami testów}
            \end{itemize}
        \end{block}
        \column{0.46\textwidth}
        \begin{exampleblock}{Rekomendacje przed wdrożeniem}
            \small
            \begin{itemize}
                \item monitoring samego CrowdSec (status LAPI, błędy bouncera),
                \item świadoma decyzja fail-open / fail-close,
                \item eskalacja czasu bana dla recydywistów (profile),
                \item whitelista tylko dla adresów administratora,
                \item obrona warstwowa: aktualizacje WP, MFA dla panelu, rate limiting w nginx, kopie zapasowe.
            \end{itemize}
        \end{exampleblock}
    \end{columns}
\end{frame}

% =====================================================================
\section{Źródła}
% =====================================================================

\begin{frame}{Bibliografia}
    \kto{OSOBA 3}
    \scriptsize
    \begin{thebibliography}{99}
        \bibitem{csdocs} CrowdSec, \emph{Documentation}, \url{https://docs.crowdsec.net/}.
        \bibitem{csappsec} CrowdSec, \emph{AppSec Component}, \url{https://docs.crowdsec.net/docs/appsec/intro}.
        \bibitem{csbouncer} CrowdSec, \emph{cs-nginx-bouncer}, \url{https://github.com/crowdsecurity/cs-nginx-bouncer}.
        \bibitem{f2b} fail2ban, \url{https://github.com/fail2ban/fail2ban}.
        \bibitem{modsec} OWASP, \emph{ModSecurity}, \url{https://github.com/owasp-modsecurity/ModSecurity}.
        \bibitem{crs} OWASP, \emph{Core Rule Set}, \url{https://coreruleset.org/}.
        \bibitem{coraza} OWASP, \emph{Coraza WAF}, \url{https://coraza.io/}.
        \bibitem{ngx} nginx, \emph{Module ngx\_http\_core\_module -- resolver}, \url{https://nginx.org/en/docs/http/ngx_http_core_module.html\#resolver}.
        \bibitem{nist} NIST, \emph{Cybersecurity Framework 2.0}, \url{https://www.nist.gov/cyberframework}.
        \bibitem{owasp} OWASP, \emph{Logging Cheat Sheet}, \url{https://cheatsheetseries.owasp.org/cheatsheets/Logging_Cheat_Sheet.html}.
    \end{thebibliography}
\end{frame}

\begin{frame}[plain]
    \centering
    {\Huge\bfseries\color{CyberBlue} Dziękujemy!}
\end{frame}

\end{document}
